\section{Conclusion}
\label{sec:conclusion}

In this paper, we identified a fundamental structural flaw in modern Multimodal Large Language Models (MLLMs): the \textbf{cross-modal granularity mismatch}. This misalignment between continuous visual features and fragmented, discrete sub-word tokens forces autoregressive generation into an unverified, \emph{open-loop decoding state}. Consequently, the statistical inertia of language priors frequently overwhelms visual constraints, catalyzing object and relational hallucinations.

To address this root cause, we introduced \textbf{CHORD} (Calibrating Hallucinations via Object-Resonant Decoding), a training-free, plug-and-play decoding framework. A parameter-free tokenizer bridge lifts fragmented tokens into semantically complete spans. CHORD then computes the \textbf{Resonance Delta ($\Delta S$)} in an external, shared semantic space, amplifying visually grounded entities and penalizing fabricated content without brittle heuristics or heavy black-box evaluators.

We further mitigated the latency bottleneck common to decode-time interventions by using \textbf{continuous Top-$1$ monitoring} with KV-cache acceleration instead of naive entropy-only gating, reducing overhead while still catching overconfident hallucinations.

Experiments across multiple MLLM families show that CHORD substantially reduces object hallucinations without sacrificing linguistic fluency. CHORD demonstrates that strong, cross-model grounding control can be achieved through external object-resonant decoding, offering a scalable paradigm for more trustworthy MLLM deployment.
